
\def\PUDcodigo{ECA215}

\def\PUDcodacad{04507.24}

\def\PUDdisciplina{Circuitos Elétricos II}

\def\PUDsemestre{S5}

\def\PUDhoras{80}

%NOVOS ---------------------------------------------------
\def\PUDhorasTeoria{64}
\def\PUDhorasPratica{16}

\def\PUDinterECA{ 
\hyperref[PUD:ECA210]{Circuitos Elétricos I (04507.20)}

\hyperref[PUD:ECA222]{Instalações Elétricas (04507.30)}

\hyperref[PUD:ECA223]{Máquinas Elétricas (04507.36)} 
}

\def\PUDatuaisECA{---}

\def\PUDrecursos{Projetor multimídia; Quadro branco e pincel; Aula em laboratório de eletricidade CA.}

%----------------------------------------------------------

\def\PUDcreditos{4}

\def\PUDprerequisito{ \hyperref[PUD:ECA210]{ECA210} }

\def\PUDpreacad{ \hyperref[PUD:ECA210]{Circuitos Elétricos I (04507.20)} }

\def\PUDobjetivos{Conhecer as grandezas elétricas envolvidas em um circuito de corrente alternada; Calcular correntes e tensões elétricas em dispositivos passivos sob corrente alternada; Analisar o comportamento das correntes e tensões utilizando fasores; Compreender como se comporta a potência em circuitos CA; Analisar o comportamento de circuitos polifásicos, em especial os trifásicos; Entender o fenômeno da ressonância e suas consequências em circuitos CA; Compreender os conceitos introdutórios sobre a transformada de Laplace.}

\def\PUDementa{Sinais de Tensão e Corrente no Domínio do Tempo e da Frequência em CA;  
Representação Fasorial;  Impedância;  Admitância;  Métodos de Análise e 
Teoremas de Circuitos;  Potência em Circuitos CA;  Circuitos Polifásicos;  
Ressonância;  Introdução à Transformada de Fourier.}

\def\PUDconteudo{ & Conceitos em Corrente Alternada: Circuitos de Corrente Alternada.  Formas de Ondas de Tensões/Correntes.  Conceitos: ciclo, período,
freqüência, velocidade ou frequência angular;  valor de pico;  fase;  defasagem;  valor médio e valor eficaz para tensões e correntes senoidais.  
Tensão Senoidal em Circuitos com Resistor, Indutor e Capacitor.  Tensão senoidal em circuitos RLC.  
\\ 
 & Senóides e Fasores: Representação fasorial de uma corrente ou tensão alternada senoidal.  Impedância.  impedância equivalente.  Diagramas
Fasoriais.  Admitância.  Métodos de resolução de circuitos CA usando impedância e admitância. 
\\ 
 & Análise Senoidal em Regime Permanente: Análise de Circuitos em CA.  Estrela-Triângulo.  Análise nodal.  Análise de malha.  Superposição. 
Transformação de Fontes. Teoremas de Thevenin e Norton.    
\\ 
 & Análise da Potência CA: Potência instantânea, média, ativa, reativa, aparente e complexa.  Medição de potência ativa.  Fator de potência. 
Correção do fator de potência. 
\\ 
 & Circuitos Polifásicos: Seqüência de fases.  Conexões das cargas em estrela e delta.  Tensões e correntes de fase e de linhas.  Diagramas
fasoriais.  Transformações Estrela-Delta e Delta-Estrela.  Métodos de circuitos equilibrados.  Potência em circuitos trifásicos.  Medição de
potência ativa e potência reativa.  Circuitos Desequilibrados. 
\\ 
 & Ressonância: Ressonância série e paralela.  Conceitos básicos de filtros. 
\\ 
 & Transformada de Laplace: conceitos e definições, propriedades, teoremas, transformada inversa de Laplace, expansão em frações parciais, resolução de equações diferenciais via transformada de Laplace, análise de circuitos elétricos via transformada de Laplace.
 %Transformada de Fourier: Introdução à Transformada de Fourier. 
 }

\def\PUDmetodo{
%Aulas expositórias que podem ser teóricas e/ou práticas, onde as práticas no laboratório serão marcadas no decorrer da disciplina.
Aulas teóricas expositivas com auxílio de recursos audio-visuais e atividades práticas em laboratório com roteiros definidos.}

\def\PUDavalia{A avaliação será feita com aplicação de provas teóricas e/ou práticas, além da possibilidade de inclusão de trabalhos e seminários no decorrer da disciplina.}

\def\PUDbibbasica{ & \href{www.biblioteca.ifce.edu.br/index.asp?codigo_sophia=62512}{Alexander}, C. K.; Sadiku, M. N. O.  Fundamentos de Circuitos Elétricos.  5º ed.  Porto Alegre, RS: AMGH, 2013.  874p.  ISBN: 9788580551723.
\\ 
 &  \href{www.biblioteca.ifce.edu.br/index.asp?codigo_sophia=41345}{Nilsson}, J. W.; Riedel, S. A.  Circuitos Elétricos.  8º ed.  São Paulo, SP: Pearson Prentice Hall, 2009.  574p.  ISBN: 9788576051596.
\\ 
 & \href{www.biblioteca.ifce.edu.br/index.asp?codigo_sophia=61748}{Boylestad}, R. L.  Introdução à Análise de Circuitos.  12º ed.  São Paulo, SP: Pearson Prentice Hall, 2012.  959p.  ISBN: 9788564574205. }

\def\PUDbibcomplementar{ & \href{www.biblioteca.ifce.edu.br/index.asp?codigo_sophia=23920}{Irwin}, J. D.  Introdução à Análise de Circuitos Elétricos.  1º ed.  Rio de Janeiro, RJ: LTC, 2005.  391p.  ISBN: 8521614322. 
\\ 
 & \href{www.biblioteca.ifce.edu.br/index.asp?codigo_sophia=23867}{Halliday}, David.  Física 3.  5º ed.  Rio de Janeiro, RJ: LTC, 2004.  377p.  ISBN: 8521613911.
\\ 
 & \href{www.biblioteca.ifce.edu.br/index.asp?codigo_sophia=24602}{O’Malley}, J.  Análise de Circuitos.  2º ed.  Rio de Janeiro: LTC, 2005.  677p.  ISBN: 8534601194.
\\
 & \href{www.biblioteca.ifce.edu.br/index.asp?codigo_sophia=55600}{Mariotto}, Paulo Antônio.  Análise de Circuitos Elétricos.  E-book.  Pearson.  390p.  ISBN: 9788587918062.
\\
 & \href{www.biblioteca.ifce.edu.br/index.asp?codigo_sophia=55510}{Burian Júnior}, Yaro; Lyra, Ana Cristina Cavalcanti.  Circuitos Elétricos.  E-book.  Pearson.  320p.  ISBN: 9788576050728. }

\def\PUDautor{José Daniel de Alencar Santos}

\def\PUDdata{2013-04-23}

\def\PUDrevisao{2}

\def\PUDrevisor{José Daniel de Alencar Santos}

\def\PUDdatarevisao{2019-05-15}

\PUDcorpo
